%%%%%%%%%%%%%%%%%%%%%%%%%%%%%%%%%%%%%%%%%%%%%%%%%%%%%%%%%%%%%%%%%%%%%%%%%%%%%%%%
%2345678901234567890123456789012345678901234567890123456789012345678901234567890
%        1         2         3         4         5         6         7         8
%% memo.tex
%% V1.0
%% 2015/09/20
%% by Prof. Rui Santos Cruz
%% based on a template created by Rob Oakes
%%%%%%%%%%%%%%%%%%%%%%%%%%%%%%%%%%%%%%%%%%%%%%%%%%%%%%%%%%%%%%%%%%%%%%%%%%%%%%%%
\documentclass[a4paper,12pt]{texMemo}
% --> Please Choose the MAIN LANGUAGE for the document in package BABEL
% --> by replacing "main=" in language name selector. Default is "main=english"
\usepackage[main=english,portuguese]{babel} % Defines Main language
\usepackage[utf8]{inputenc}
\usepackage{iflang}
\usepackage{ifthen}
\usepackage{parskip}
\setlength{\parindent}{15pt}
\logo{\includegraphics[width=0.3\textwidth]{madison.jpg}} % Logo
%%%%%%%%%%%%%%%%%%%%%%%%%%%%%%%%%%%%%%%%%%%%%%%%%%%%%%%%%%%%%%%%%%%%%%%%%%%%%%%%
%	MEMO INFORMATION --> Write your info in the following tags.
%%%%%%%%%%%%%%%%%%%%%%%%%%%%%%%%%%%%%%%%%%%%%%%%%%%%%%%%%%%%%%%%%%%%%%%%%%%%%%%%
\memofrom{Luxi Cao, Zongyan Wang} % Sender(s) Name
\memocourse{Stat 992: Dependence Modeling with Copulas} % Course Name, or abbreviated acronym
\memosubject{Lecture 7 Scribe} % Subject
\memodate{\today} % Date, -> set to \today for automatically print todays date
%%%%%%%%%%%%%%%%%%%%%%%%%%%%%%%%%%%%%%%%%%%%%%%%%%%%%%%%%%%%%%%%%%%%%%%%%%%%%%%%
\begin{document}
\maketitle % Print the memo header information
%%%%%%%%%%%%%%%%%%%%%%%%%%%%%%%%%%%%%%%%%%%%%%%%%%%%%%%%%%%%%%%%%%%%%%%%%%%%%%%%
%	MEMO CONTENT --> Your content is written here
%%%%%%%%%%%%%%%%%%%%%%%%%%%%%%%%%%%%%%%%%%%%%%%%%%%%%%%%%%%%%%%%%%%%%%%%%%%%%%%%
Q1. How to relate infinitely divisible, domain of attraction, stable, slowly varying?
\begin{enumerate}
\item In case of multivariate distribution F, as long as $F^q, q>0$, F is infinitely divisible. if q is a positive integer, then it is always true. Otherwise, you have to check the power function.
In probability theory, a probability distribution is infinitely divisible if it can be expressed as the probability distribution of the sum of an arbitrary number of independent and identically distributed random variables. 
\item Domain of attraction: Limit distribution of Minimum or Maximum of random variables exists
\item Sum stable: For any distribution, if a random variable can be expressed as a linear combination of independent variables. Those random variables follow the same distribution family.
\item Max stable: take the maximum of a number of random variables. The random variable has a limit distribution with indivisible property. Any of these multivariate extreme value distribution can be expressed as a linear combination of unit $Fr\acute{e}chet$ random variables.
\item slowly varying: many of those joint survival distribution can be expressed as $\bar{F}(x) \sim x^{-\alpha}\dot L(x)$
\end{enumerate}

Q2. Tow-stage estimation for copula models, why a big difference between $\hat{\eta}$ and $\tilde{\eta}$, relative to the standard errors for V, suggests either the copula model or the uni-variate models may be inadequate. How 'big' is a relative big?
\begin{enumerate}
\item If our estimation procedure is consistent then both good enough. This is so called consistency check: $\hat{\eta}$ and $\tilde{\eta}$ should not differ a lot.
\item Some people dislike 2 stage, because after you estimate the first part, there is a random variable. Then you have to adopt return part. Then the second part automatically isn't good.
\item Coordinate decent is more popular recently, estimate the parameter one by one. 
\item How big is a relative big? Ex. in linear regression model, how to consider two variances are the same? Rule of Thumb: $\hat{\eta}$/$\tilde{\eta}$ $\in [0.5, 2]$
\end{enumerate}

Q3. How to use Laplace transformation to construct a copula can be further studied. Chap3 would talk about it.

4. Using frailty parameter and resilience parameter, we can have Archimedean copulas. What is the importance of Archimedean copula?

Archimedean Copula's property will be shown. Archimedean Copula is a very broad family. Many of copulas dependence structures could be constructed with the help of Archimedean Copula. it generates Gumbel Copulas which is extreme value copula.

5. Is there any formal definition of concordance ordering? since the concordance ordering is a way to compare multivariate distributions with the same set of uni-variate margins, but usually what we know is joint distribution rather than marginal distribution. So how do we apply this to practice problem?

\begin{enumerate}
\item The first part is discussed in Section 2.9
\item Second part: Partly agree with what we know is joint distribution rather than marginal distribution. If we consider Multi-Normal, then it is the case. 
\item The reason for we study Copulas dependence is that we don't have unified parametric form to model everything. Since in many cases, those marginal distributions are not within the same family. 
\item For those marginal distribution, we always could do some transformations. Firstly, we fit some distribution to the marginal, then transform to uniform or $Fr\acute(e)chet$, then we use joint structure to model the joint distribution. This is the sample statement.
\item Applications will be shown
\end{enumerate}

6. In the slides, there is one paragraph said that: 
"The concordance ordering is a way to compare multivariate distributions with the same set of univariate margins so that we can decide if one multivariate distribution represents more dependence that another."
I wonder how to use the concordance ordering?

It will be taught soon. Be patient little little baby.

text p28: "2.1.2 Absolutely continuous and singular components". It says: "The decomposition could also be written in terms of survival functions." Could you please explain more about this?

Note $F = \delta_{ac}F_{ac} + \delta_{sin}F_{sin} + \delta_{d}F_d$ could be written in terms of survival function:
$1 - F = \bar{\delta}_{ac}(1 - F_{ac}) + \bar{\delta}_{sin}(1 - F_{sin}) + \bar{\delta}_{d}(1 - F_d)$. Those three delta are different.

8.VC will later be talked about
estimate PC, hold off


9. apply C, eg:
$Y=(y_1,y_2,....,y_n)$    vector $\lambda(x_1,x_2,x_3,..,x_p)=f(a_1x_1+a_2x_2+...a_px_p)$
multivariavte normal choose lambda as mu would be the same thing.

10. compare C and Multivariate dist fn
MD can also be considered as C, but it doesn't has uniform dist
After we have transformed to the same distribution, 

11. l(tx)/l(x) has a limit 1. eg. logx, 

12. 3 distinct regions:
In general linear regression, many situations we always have some good model fitting. 
Ratio estimation procedure, good estimation at upp tail.
No concrete answer.
Extreme dependence no much attention to the middle part.
Regression pay attention to the middle part, which dominates the model. 
We don't use residual, we check misclassification rate. 
If you generate data using the hyperplane,. either solution is acceptable for classification. 

13. VC in 2.6 chp4 talk more

14. high order moment
C trandform MV to uniform, the high order moment exists.

15. CC in a subset we choose C. Up to our purpose.

16. truncation 
orig: model fitting, truncation, 
T,C applies to univariate marginals 
tail dependence: $P(X>u|Y>u)$
Bivariate normal,  $lim_{u->u^*}P(X^*>u|Y^*>u)=\lambda?0$
$X^*=min{X,u^*},Y^*=min{Y,u^*}$   

17. Effect of misspecification
Truncation and censoring effect: Research question more like survival analysis part.
Basically, truncation adn censoring are different. If we treat truncated data as normal distribution and transform it to uniform distribution, then we don't know what will happen yet. It also depends on the types of model we have. It's good research question. 
18.frechet class:
eg:copula fn 
first one. 
Upper bound:always is a joint distribution, lower bound: not always a joint distribution

Q3.
%%%%%%%%%%%%%%%%%%%%%%%%%%%%%%%%%%%%%%%%%%%%%%%%%%%%%%%%%%%%%%%%%%%%%%%%%%%%%%%%
\end{document}